\scanpage{105}
\[
\mu:\underbrace{\bigwedge^2 B/k}_{\substack{\parallel\\\det B/k}}\longrightarrow A/B
\]

Or on a un iso. can.
\margin{attention à conv. de signes}
\[
\det B\simeq\underbrace{\det k}_{k}\otimes\det B/k\simeq\det B/k
\]
donc on peut interpréter $\mu$ comme
\[
\mu:\underset{\substack{\rotatebox{90}{$\simeq$}\\(A/B)^{\otimes-1}}}{\det B}\longrightarrow A/B
\qquad\deleted{(\unsure{\simeq}\det B^{\unsure{-\otimes1|-\otimes}})}
\]
soit encore
\[
\mu\in A/B^{\otimes2}\simeq D^{-\otimes2}
\]

Ceci dit, il est clair que
\[
B\text{ sous-algèbre}\Longleftrightarrow\mu=0
\]
Il suffit de voir que
\[
\mu=\pm\delta
\]
\margin{le signe dépendant des conventions de signes}

Or il suffit de le voir dans le cas universel $S=\mathbb{P}^2_{\mathbb{Z}}$, et comme \unsure{nos|mes|ces} \deleted{il n'y a pas d'éléts} $\mu$ et $\delta$ \deleted{sont les $=\text{s}$} universels ont les $=\text{s}$ \deleted{zéros, \unsure{sans}} ens de zéros, on gagne !
% archive sheet number: 104
